\section{Trabajo futuro}
\label{sec:concl_futuro}
% Extensión estimada: 1 página, una oración de entrada, una lista de cinco líneas y un párrafo de cierre (ficha 7.4).
% Fuentes: sec:disc_limitaciones (alcance, ancla de una ejecución, decisiones sin evaluar), sec:disc_concordancia,
% sec:disc_texto_imagen, sec:disc_limites. Sin futurología (guía §3): cada línea dice qué falta evaluar, no qué
% beneficio se espera.

Las líneas de trabajo futuro se desprenden de las limitaciones de la Sección~\ref{sec:disc_limitaciones}, y cada una
consiste en evaluar algo que este trabajo dejó sin medir:
\begin{itemize}
    \item \textbf{La evaluación por especialistas.} La continuación más directa es que dermatólogos lean los reportes
    junto con sus imágenes, para verificar que los rasgos descritos estén en la lesión y juzgar si el reporte sería
    aceptable en la práctica. Esa evaluación, junto con un conjunto de reportes de referencia redactados por
    especialistas para las imágenes dermatoscópicas, permitiría pasar de la fidelidad diagnóstica a la validez
    clínica.
    \item \textbf{Una regla de decisión orientada al costo clínico.} El ancla del \emph{pipeline} usa la clase más
    probable de Swin-Tiny, no el umbral de melanoma, y el \emph{pipeline} se eligió por una métrica que trata igual
    todos los errores (Sección~\ref{sec:disc_concordancia}). Queda por evaluar cómo cambian la consistencia por clase y
    la proporción de casos derivados si el ancla usa ese umbral, o si el ruteo pondera los errores según su costo.
    \item \textbf{Otras señales de confianza.} Los errores del clasificador y del \gls{knn} coinciden unas tres veces
    más de lo que coincidirían por azar, posiblemente porque los dos aprenden de las mismas imágenes. Queda por
    evaluar si la probabilidad del clasificador, que está bien calibrada, o un conjunto de casos de referencia
    distinto del de entrenamiento detectan una parte de esos errores compartidos.
    \item \textbf{La validación externa del \emph{pipeline}.} En \emph{HIBA} se evaluó solo el clasificador. Aplicar
    el \emph{pipeline} completo a otra población, y a otras modalidades como la fotografía clínica, permitiría medir la
    fidelidad del reporte fuera de \emph{HAM10000} en lugar de inferirla.
    \item \textbf{Otros anclajes y generadores.} La relación entre la exactitud del clasificador y la fidelidad del
    reporte se caracterizó con tres fuentes y una sola ejecución de Swin-Tiny. Un barrido de clasificadores y de
    ejecuciones la describiría con más detalle. Del mismo modo, quedan por evaluar los generadores que razonan antes de
    diagnosticar y los de mayor tamaño, junto con una medida de qué parte de la descripción proviene de la imagen y
    qué parte del diagnóstico recibido (Sección~\ref{sec:disc_texto_imagen}).
\end{itemize}

Todas estas líneas parten de la misma conclusión: en un sistema que combina un clasificador y un generador, la
fidelidad diagnóstica del reporte depende del clasificador y de la señal que decide cuánto afirmar. Por eso las
mejoras que este trabajo deja planteadas apuntan a esos dos componentes y a su evaluación con especialistas, más que al
tamaño del generador.
